\documentclass[10pt,a4paper]{amsart}
\usepackage[margin=19mm]{geometry}
\usepackage{amsmath,amssymb,mathtools}
\usepackage[scr=rsfs,cal=euler]{mathalfa}
\usepackage{xfrac}
\usepackage[unicode,hidelinks,bookmarks=false]{hyperref}
\newcommand{\ie}{i.e.\ }
\newcommand{\cf}{cf.\ }
\newcommand{\resp}{respectively\ }
\newcommand{\Subsection}{\subsection}
\providecommand{\nobreakdash}{\nobreak\hbox{-}}
\setlength{\emergencystretch}{2em}
\multlinegap0pt
\usepackage{bm}
\usepackage{bbm}
\usepackage{dsfont}
\usepackage{xspace}
\usepackage{cancel}
\usepackage{mathtools}
\usepackage{amsthm}
\makeatletter
\@ifundefined{lemma}{\newtheorem{lemma}{Lemma}}{}
\makeatother
\DeclareMathOperator{\vol}{vol}
\newcommand{\cG}{\mathcal{G}}
\newcommand{\loy}{\mathcal{L}}
\title[Approximating temporal modularity on graphs of small underly]{Approximating temporal modularity on graphs of small underlying treewidth}
\author{Vilhelm Agdur, Jessica Enright, Laura Larios-Jones, Kitty Meeks, Fiona Skerman, Ella Yates}
\date{Source: arXiv:2507.17541v1 (CC BY 4.0)}
\begin{document}
\maketitle

\noindent\emph{Source and licence.} Approximating temporal modularity on graphs of small underlying treewidth. Version arXiv:2507.17541v1 (CC BY 4.0), distributed under the Creative Commons Attribution 4.0 International licence, as recorded in the arXiv licence metadata for this exact version and shown on the abstract page https://arxiv.org/abs/2507.17541v1. Full rights notice: licenses/arxiv-2507.17541-CC-BY-4.0.txt.\par\smallskip

\noindent\textit{Editorial scope.} Counted unit: the Lemma whose source label is \texttt{lem:tw root} in the article (source0.tex lines 738--741, source label \texttt{lem:tw root}), delivered together with the article's own complete original proof of it (source0.tex lines 742--782, one \texttt{proof} environment of 41 lines). No mathematics is written, repaired or re-derived: statement and proof are the article's own text line for line. Required context retained uncounted: none. Every definition, display and label of those lines is retained unchanged. Omitted from this package and not redistributed: 1--737, 783--892. Nothing inside a retained excerpt is omitted or altered: every retained body is delivered exactly as the article's lines are written, so no omission receipt of any kind accompanies this package. Citations: the retained excerpts carry no citation command at all, so no bibliography entry of the article is transcribed and no reference is redistributed. Reference adaptation: none was needed and none was made; every reference inside the delivered unit resolves inside it, so no unresolved reference or citation is produced. Printed slips kept literal: no printed word, symbol, hypothesis or number of the counted statement or of its proof was altered. Layout and class: the article's own class is \texttt{llncs} with packages including fontenc, graphicx, inputenc, apxproof, amsthm, amsmath, amssymb, latexsym, color, soul, mathtools, hyperref, cleveref, framed; the wrapper is the lane's pinned \texttt{amsart} editorial wrapper at 10pt on A4, which restates the theorem-like environments the retained text needs and adds the article's own macro definitions that the retained text needs (\texttt{\textbackslash cG}, \texttt{\textbackslash loy}, \texttt{\textbackslash vol}), with extra wrapper packages (bm, bbm, dsfont, xspace, cancel, mathtools). The delivered numbering is the unit's own. Assets: no retained span contains a figure environment, a graphics-inclusion command, a PDF-page inclusion, a table or a TikZ picture; no image file is copied into this package, the package declares no asset dependency and no image file is redistributed with it. Licence: version arXiv:2507.17541v1 of this article is distributed under the Creative Commons Attribution 4.0 International licence, as recorded in the arXiv licence metadata for this exact version and shown on the abstract page https://arxiv.org/abs/2507.17541v1. A full rights notice is delivered at licenses/arxiv-2507.17541-CC-BY-4.0.txt. Characters: every retained excerpt is pure ASCII, so the delivered engine, XeLaTeX with its default Latin Modern text font, reports no missing character.
\begin{lemma}\label{lem:tw root}
Let $\cG$ be a temporal graph with $n$ vertices, $m$ edges and lifetime $T$.
Given the set of all valid states at the root of a nice tree decomposition of $\cG$, the maximum non-normalised temporal $c$-modularity of $\cG$, $\widetilde{q}^{\,*}_c(\cG)$, can be found in time $O(2^{cT}m^{c(T+1)}T^{c+2}cn)$. 
\end{lemma}

\begin{proof}
We claim that
\begin{equation}\label{eqn:root_lemma}
\widetilde{q}^{\,*}_c(\cG)=\max_{\text{valid state at $b$ }(\pi,\alpha,\beta,\gamma)} \left(\sum_{p \in [c]} \left( 2\alpha(p) -\sum_{t=1}^T  \frac{\beta(p,t)^2}{2m_t}\right) + \omega \gamma\right).
\end{equation}
First we want to show that 
\[\widetilde{q}^{\,*}_c(\cG)\geq\max_{\text{valid state at $b$ }(\pi,\alpha,\beta,\gamma)} \left(\sum_{p \in [c]} \left( 2\alpha(p) -\sum_{t=1}^T  \frac{\beta(p,t)^2}{2m_t}\right) + \omega \gamma\right).\]
To do this we consider the valid state $(\pi,\alpha,\beta,\gamma)$ where the maximum is obtained. As the state is valid, it is supported by a partition function $\pi^*$ on all forgotten vertices. We know at the empty root node that all vertices have been forgotten, hence $\pi^*$ is a partition function on all vertices on the graph. 
We show that the sum we maximise in Equation \ref{eqn:root_lemma} is therefore the modularity of the partition given by $\pi^*$. 
By validity $\alpha$ exactly counts the edges within each part. This means for each part $p$ we have that $\alpha(p)=\sum_{t=1}^Te_{G_t}(p)$. Therefore, we have $\sum_{p \in [c]}2\alpha(p)=\sum_{p\in[c]} 2\sum_{t=1}^T e_{G_t}(p)=\sum_{t=1}^T \sum_{p \in [c]} 2e_{G_t}(p))$.
    We know from validity that $\beta$ exactly counts the degree in each part and time. This means for any part $p$ and time $t$ $\beta(p,t)=\vol_{G_t}(p)$.
    Similarly $\gamma$ is exactly the loyalty count, so $\gamma=\loy(\pi^*)$. 
\begin{align*}
\left(\sum_{p \in [c]} \left( 2\alpha(p) -\sum_{t=1}^T  \frac{\beta(p,t)^2}{2m_t}\right) + \omega \gamma\right)\\
 =\sum_{t=1}^T \sum_{p \in [c]} \left( 2e_{G_t}(p) - \frac{\text{vol}_{G_t}(p)^2}{2m_t}\right) + \omega \loy(\pi^*)\\
=\widetilde{q_c}(\cG,\pi^*) \leq \widetilde{q}^{\,*}_c.
\end{align*}

Next we show that
\[\widetilde{q}^{\,*}_c(\cG)\leq\max_{\text{valid state at $b$ }(\pi,\alpha,\beta,\gamma)} \left(\sum_{p \in [c]} \left( 2\alpha(p) -\sum_{t=1}^T  \frac{\beta(p,t)^2}{2m_t}\right) + \omega \gamma\right)\]
We know that at each step every possible combination of functions a state could take is considered and accepted as valid if it is supported by some partition function $\pi^*$ on the forgotten vertices. This means that, at a given node, every possible partition function on the forgotten vertices has a corresponding valid state. Therefore, at the root node where all vertices are forgotten every possible partition function on the temporal graph has a corresponding valid state.

Consider a valid state at the root $(\pi,\alpha,\beta,\gamma)$.
With a supporting partition function $\pi^*$ acting on all the vertices of the graph.
We start from the definition of non-normalised temporal $c$-modularity on the partition $\pi^*$:
    \[
        \widetilde{q}_\omega(\cG,\pi^*) = \sum_{t=1}^T \sum_{p \in [c]} \left( 2e_{G_t}(p) - \frac{\text{vol}_{G_t}(p)^2}{2m_t}\right) + \omega \loy(\pi^*)
    \]
    By validity $\alpha$ exactly counts the edges within each part. This means for each part $p$ we have that $\alpha(p)=\sum_{t=1}^Te_{G_t}(p)$. Therefore, we have $\sum_{t=1}^T \sum_{p \in [c]} 2e_{G_t}(p))=\sum_{p\in[c]} 2\sum_{t=1}^T e_{G_t}(p)$$=\sum_{p \in [c]}2\alpha(p)$.
    We know from validity that $\beta$ exactly counts the degree in each part and time. This means for any part $p$ and time $t$ $\beta(p,t)=\vol_{G_t}(p)$.
    Similarly $\gamma$ is exactly the loyalty count, so $\gamma=\loy(\pi^*)$.
    This gives non-normalised temporal modularity of the partition function $\pi^*$ as:
    \[
        \widetilde{q}_\omega(\cG,\pi^*) = \sum_{p \in [c]} \left( 2\alpha(p) -\sum_{t=1}^T  \frac{\beta(p,t)^2}{2m_t}\right) + \omega \gamma
    \]
The non-normalised temporal modularity of a partition can be clearly calculated given a valid state, as we are considering all possible partitions this will include the maximum which is equal to $\widetilde{q}^{\,*}_c(\cG)$.
    
For each state the non-normalised temporal modularity of a supporting partition can be calculated in time $\mathcal{O}(cT)$.
To find the non-normalised temporal $c$-modularity of the graph we compare the calculated modularity of all root states. All the vertices have been forgotten so the partition function $\pi$ of the valid state is always empty, therefore there are $(mT+1)^c(2m+1)^{cT}(nT+1)$ possible valid states. Hence at the root the maximum non-normalised temporal $c$-modularity can be found in $\mathcal{O}(2^{cT}m^{c(T+1)}T^{c+2}cn)$.
\qed
\end{proof}

\end{document}
